从某种意义上说，语句集合描述了它们在其中共同为真的结构；这些结构就是它们的\emph{模型}。反过来，如果我们从一个结构或一组结构出发，可能会对以它们为模型的语句集合感兴趣，这个集合就是该结构或该组结构的\emph{理论}。任何这样的语句集合都具有如下性质：被它们所语义后承的每一个语句都已经在该集合中；它们是\emph{封闭的}。更一般地，如果集合 $\Gamma$ 在语义后承下封闭，我们就称 $\Gamma$ 为一个理论；并且，如果 $\Gamma$ 恰好由 $\Delta$ 所语义后承的所有语句组成，我们就说 $\Gamma$ 被~$\Delta$ \emph{公理化}了。

数学提供了许多理论的例子，例如线性序理论、群论，或者算术理论——例如由 Peano（皮亚诺）公理公理化的理论。但其他学科中也有许多重要理论的例子，例如，关系数据库可以被看作理论，形而上学则关注可以被公理化的部分论理论。

在建立一个理论加以研究时，一个重要问题是它的语言是否具有足够的表达力，让我们能表达希望该理论讨论的一切；另一个问题是它是否足够强，能证明我们希望它证明的结论。要\emph{表达}一个关系，我们需要一个含有相应数量自由变元的公式。在\emph{集合论}中，我们只有 $\in$ 作为谓词符号，但它允许我们用$\lforall[u][(u \in x \lif u \in y)]$来表达 $x \subseteq y$。事实上，\emph{Zermelo-Fraenkel（策梅洛-弗兰克尔）集合论}~$\Log{ZFC}$ 足够强，既能（几乎）表达每一个数学论断，又能（几乎）证明每一个数学定理，而它仅仅使用了少数几条公理和一系列日益复杂的定义，比如~$\subseteq$ 的定义。